\section{Standard forms and finite completion}\label{sec:structure}
\subsection{Descent of the mesh presentation}

We first explain why geometric standardness gives a presentation over the original field. For a finite-dimensional algebra $C$, its radical grading is
\[
 \gr C=\bigoplus_{j\geq0}\rad^j C/\rad^{j+1}C.
\]
A radical-graded algebra is an algebra isomorphic to this graded algebra, with its grading forgotten.

\begin{lemma}\label{lem:grading}
Let $k$ be a finite field. If $C_{\bar k}$ is radical-graded, then $C$ is radical-graded.
\end{lemma}
\begin{proof}
Put $B=\gr C$. Scalar extension commutes with the radical because $k$ is perfect. Choose a geometric isomorphism $C_{\bar k}\to B_{\bar k}$. Its associated graded map is a graded automorphism of $B_{\bar k}$; composing with its inverse makes the associated graded map the identity.

Let $U$ be the reduced algebraic group of automorphisms of $B$ inducing the identity on its associated graded algebra. This group is smooth since $k$ is perfect. It is connected: write an automorphism in homogeneous blocks. Its diagonal blocks are identities and its other blocks go from degree $a$ to a strictly larger degree $b$. Conjugation by the grading action multiplies these blocks by $t^{b-a}$. It extends to $t=0$, with value the identity. Thus every geometric point of $U$ belongs to the identity component. The isomorphism scheme of maps $C\to B$ inducing the graded identity is geometrically nonempty. Choosing one geometric isomorphism identifies it, by composition, with the automorphism scheme inducing the graded identity. Reduction commutes with extension of the perfect field $k$, so its reduced scheme becomes $U_{\bar k}$ with its translation action. It is therefore a smooth $U$-torsor over $k$. Lang's theorem \cite{Lang} gives a rational point of this torsor.
\end{proof}


Take one representative of each indecomposable $A$-module and let $G$ be their direct sum. Over $\bar k$, $G_{\bar k}$ contains every geometric indecomposable. Indeed, any geometric indecomposable is defined over a finite separable extension $E/k$, and is a summand of the scalar extension of its restriction to $k$. Also
\[
 \End_A(G)/\rad\End_A(G)=\prod_i\FF_{q^{d_i}},
\]
so after splitting the fields its primitive summands have multiplicity one. Thus $\End_A(G)_{\bar k}$ is the basic endomorphism algebra of all geometric indecomposables.

In a mesh category, arrows are irreducible morphisms and the defining relations are homogeneous of degree two. In finite representation type its arrow ideal is nilpotent, so the grading is the radical grading. Lemma~\ref{lem:grading} descends this grading to $\End_A(G)$. One can align the vertex idempotents with the original summand projections: if two orthogonal systems $(e_i)$ and $(f_i)$ coincide modulo the radical, then
$u=\sum_i f_i e_i$ lies in $1+\rad\End_A(G)$ and satisfies $f_i u=u e_i$.

Here is the resulting modulated presentation. A partial translation quiver has a partially defined injective translation $\tau$ on vertices. For every arrow $a:y\to x$ with $\tau x$ defined there is a companion
\[
 \sigma a:\tau x\longrightarrow y;
\]
There is also an inverse companion $\sigma^{-1}a:x\to\tau^{-1}y$ whenever $\tau^{-1}y$ is defined. These partial operations are mutually inverse, and $\sigma^2$ is translation on arrows wherever defined. The companion pairings are nondegenerate. These assertions hold in the Auslander--Reiten quiver and remain true after deleting projectives whenever all the indicated endpoints remain. Over a splitting field, the relation at $x$ is the sum of the length-two paths $a\,\sigma a$, with the coefficients from its almost-split sequence. Over $k$, the vertices have their residue fields, the arrow spaces are the corresponding irreducible-morphism bimodules, and the relations are the coevaluation tensors for the almost-split pairings. The geometric quadratic presentation implies the presentation over $k$: scalar extension is faithfully flat, so the descended degree-two relations generate the entire kernel.

The stable Auslander algebra
\[
 \Lambda=\End_A(G)/[\operatorname{proj} A]
\]
is obtained by deleting the projective vertices and every path through them. A map between nonprojective indecomposables factoring through a projective lies in the square of the radical, since neither factor can split. Thus this quotient does not change the remaining arrow spaces. Its relation at an old vertex is exactly the sum of the remaining length-two mesh paths. The sum can be empty. Translation is retained only when both translated endpoints remain.

\subsection{Periodic components}

Call a vertex periodic if every positive and negative iterate of its translation exists. The full subquiver on periodic vertices of a representation-finite Auslander--Reiten quiver is a disjoint union of quotients of repetitive Dynkin quivers; this is the stable-subquiver statement of Riedtmann's structure theorem \cite[Hauptsatz]{Riedtmann}. It applies to the stable full subquiver of an arbitrary representation-finite algebra. Its vertices need not constitute all nonprojective vertices.

We record the cuts and their finite-field forms needed for counting. For Dynkin trees other than $A_{2n}$, orient the tree away from its central fixed vertex, or choose an orientation invariant under its diagram automorphisms. The automorphisms of the repetitive quiver are generated by translation and diagram automorphisms preserving such a cut. For $A_{2n}$ there is instead a generator $\rho$ with $\rho^2=\tau$. These assertions, including the exceptional generator, are given explicitly in \cite[Proposition 10]{Bongartz23}. 

In the first case, the forward arrows form a finite union of Dynkin species $H$; the companions are the transpose arrows. Deck transformations and finite-field descent preserve this division. The automorphism transporting one mesh to the next becomes the species automorphism $F$ of Section~\ref{sec:counting}. Folding a positive definite Dynkin form remains positive definite: the folded form is its restriction to the invariant subspace, with the vertex-field degrees as symmetrizing weights. Lemma~\ref{lem:residues} and \eqref{eq:ordinarytotal} therefore give the total mass on every vertex-conservative extension.

The exceptional cut has the following coordinates. Label the two arms of $A_{2n}$ by $(i,p)$, where $0\leq i<n$ is distance from the centre and $p=0,1$ distinguishes the arms. If $l\in\ZZ$ is the repetition coordinate, set $t=2l+p$. Then $\tau$ changes $t$ by two and $\rho$ changes it by one, after choosing the direction of indexing. The outward arm arrows are $a_{i,t}$; the central arrows are $\varepsilon_t$. Their companions are the reverse arm arrows $b_{i,t}$. Replacing the reverse arrows at distance $i$ by $(-1)^{i+1}$ times themselves gives precisely the square-boundary equations of Section~\ref{sec:square}. A quotient by a power of $\rho$ makes $t$ cyclic. The equality of the residue-field degrees follows from the action on these coordinates. Suppose Frobenius permutes $r$ geometric components transitively, and write one as $\mathbb Z A_{2n}/\langle\rho^s\rangle$. Its $r$th power lifts to a power $\rho^a$, by the automorphism description above. At every distance $i$ from the centre it acts on $t\in\mathbb Z/s\mathbb Z$ by $t\mapsto t+a$. Every geometric vertex therefore has Frobenius orbit length $rs/\gcd(s,a)$. This length is the degree of its descended residue field, independently of $i$.

The same description holds without splitting the field. Along an arm, choose successive identifications of the vertex fields and trivializations of the invertible arrow bimodules. There is no obstruction, since an arm is a tree. At the centre, retain the invertible bimodule of each $\varepsilon_t$ and its field transport. The two summands in the central relation are the tensor product of two successive central bimodules and the first arm-return bimodule. Each projection of the relation bimodule is an isomorphism: it is one after extension to $\bar k$, hence one before extension. The internal relations identify consecutive arm-return bimodules in the same way. Choose the reverse-arrow identifications successively along the arm. This gives the displayed square-boundary equations over the original fields, with all remaining transport in the $\varepsilon_t$. Proposition~\ref{prop:squaretotal} applies to the resulting semilinear central cycle.

\subsection{Completion of the finite translation intervals}

We complete the finite partial quiver underlying $\Lambda$ by adding vertices.

\begin{lemma}\label{lem:completion}
A finite modulated partial mesh quiver of the preceding form is a full vertex quotient of a finite transpose mesh with square-boundary components. Its original periodic subquiver is unchanged. Every added translation cycle has a new vertex.
\end{lemma}
\begin{proof}
Every component of the partial permutation $\tau$ is a cycle or an interval. Give an interval coordinates $0,\ldots,l_i$, with $\tau$ decreasing the coordinate by one. Let $L=\max_i l_i$, taking $L=0$ if there are no intervals. For each old translation cycle $c$, let $m_c$ be its length and let $\alpha_c$ be the automorphism of its starting vertex field obtained by composing the prescribed field transports once around the cycle. Write $h_c$ for the finite order of $\alpha_c$. Choose an even integer $M>2L+2$ divisible by every $m_ch_c$. Then transport through $M$ steps is the identity on each old cycle. Complete each interval to a cycle of length $M$.

Follow companion arrows. A chain that is already closed has both endpoints on old periodic cycles and is left unchanged. Otherwise, if one arrow is
$a_0:(i,p)\to(j,q)$, its successive companions have coordinates
\[
 a_{2r}:(i,p-r)\to(j,q-r),\qquad
 a_{2r+1}:(j,q-r-1)\to(i,p-r).
\]
Complete each maximal open companion chain separately modulo $2M$, identifying the coordinates of an old periodic endpoint modulo its old period. Retain the chain label, so arrows from different chains are not identified merely because they have the same endpoints. If both endpoints belong to finite intervals, any new arrow with both endpoints old would require, in the even case,
$u-v=p-q+kM$. The difference between the two sides with $k=0$ has absolute value at most $2L$; hence $k=0$. For odd arrows the analogous bound is $2L+1$. Before reduction modulo $M$, the even terms with old endpoints have
\[
 \max(p-l_i,q-l_j)\le r\le\min(p,q),
\]
and the odd terms have
\[
 \max(p-l_i,q-l_j-1)\le r\le\min(p,q-1).
\]
Each endpoint of the odd range equals the corresponding endpoint of the even range or that endpoint minus one. Thus the allowed even and odd indices interlace to form one consecutive interval of companion indices. Forward and inverse companion completeness puts every term of this interval in the original chain. Thus no new old-to-old arrow occurs. With one periodic endpoint, the finite endpoint determines the unique old occurrence, while divisibility of $M$ by the old period preserves the other phase. With two periodic endpoints the chain was already closed.

Transport the vertex fields along the completed cycles. For a completed interval, choose the transport at one new closing edge to be the inverse of the product of all other transports. Its full field transport is then the identity; for old cycles this holds after $M$ steps by the choice of $M$. Extend each arrow bimodule by alternating duality relative to its two endpoint fields, using the original nondegenerate pairings wherever prescribed. After $2M$ companions the endpoint field actions agree with the original ones. Choose the last new pairing to close the chain and retain the resulting invertible linear return map as part of the species transport $F$. No identity condition on this linear return map is needed.

On each open companion chain, alternate the forward and reverse designation. Any descent automorphism preserves its two ends and commutes with taking companions, so it preserves this designation; a shift along a finite chain cannot preserve its ends. On the old closed chains use the Dynkin cuts, including the central exceptional cut, just constructed. The double-dual identifications give a species transport $F$. Choose the newly added mesh pairings to be the transpose pairings. On an old open chain, the prescribed pairings determine these choices successively. In split one-dimensional arrow coordinates, if the terms at the consecutive meshes have nonzero coefficients $c_{2r}$ and $c_{2r+1}$, take the forward arrow to be $a_{2r}$, replace the reverse arrow $a_{2r+1}$ by $c_{2r}a_{2r+1}$, and set $F(a_{2r})=-c_{2r+1}c_{2r}^{-1}a_{2r+2}$. The transpose terms then have the original coefficients at both meshes. For general bimodules the same construction uses the two duality isomorphisms in place of the scalar ratio. An open chain imposes no return condition on these choices. On the closed components the preceding Dynkin-coordinate calculation supplies the compatible choice. This gives the required transpose relations and the exceptional square relations.

Finally quotient by all new vertex idempotents. Each relation at a new vertex vanishes, since its endpoint is new. At an old vertex whose predecessor was missing, the relation has a new translated endpoint and vanishes as well. Every other old relation loses exactly the terms with new intermediate vertices. It is therefore precisely its original partial mesh relation. The image of an ideal under a quotient map is generated by the images of its generators. Thus all relations in the vertex quotient are generated by these images of mesh relations, and the original and quotient algebras have identical presentations. An added cycle contains a new vertex because $M>L+1$.
\end{proof}

The construction also applies to the opposite category. Reverse arrows and translation, and replace each bimodule by its opposite with the interchanged endpoint actions. Pairings and the square-boundary equations are carried to the corresponding opposite equations.

\subsection{Hall polynomials in the standard case}

\begin{proof}[Proof of Theorem~\ref{thm:standard}]
Apply the preceding presentation and completion to the opposite of $\Lambda$, so that its covariant representations are the right defect modules used in Section~\ref{sec:defects}. A representation of the quotient is exactly a representation of the completion with zero spaces at all new vertices. Each added translation cycle contains such a vertex, and hence cannot support a nonzero local band. Thus the full cycles supporting bands are exactly the original periodic components.

Fix a conservative extension $E/k$ of $A$. Its degree is coprime to every residue-field degree $d_i$ from Section~\ref{sec:defects}. All vertex fields of the presentation and completion are copies of these fields. In the ordinary periodic components, Lemma~\ref{lem:residues} shows that no extra extension condition is required to preserve the indecomposables of the Dynkin species. Equation~\eqref{eq:ordinarytotal} gives their total masses in $\R$. On an exceptional component, Proposition~\ref{prop:squaretotal} gives the same conclusion; replacing its field-size variable by $Q^d$ preserves $\R$.

Proposition~\ref{prop:recovery} now gives all dimension, arrow-top and arrow-socle masses supported on the original vertices. On $\Lambda$ the arrow ideal is the radical, since the mesh grading has nilpotent positive part and semisimple degree zero. These are therefore the actual module top and socle. At vertex $i$, a defect dimension $\delta_i$ measured over $k$ is dimension $\delta_i/d_i$ over its residue field; a nonintegral quotient means that the corresponding mass is zero. We have obtained exactly the masses required by the defect criterion, on every conservative extension of $A$. Corollary~\ref{cor:criterion} proves the theorem.
\end{proof}

\subsection{Ray-standard forms}
Over an algebraically closed field, a basic algebra is \emph{distributive} when its lattice of two-sided ideals is distributive. Its ray category has primitive idempotents as objects and, in each corner, the orbits under the two corner unit groups as rays. For composable morphisms $a:x\to y$ and $b:y\to z$, the ray product is zero when $ba\in b\rad A(y,y)a$, and otherwise is the ray of $ba$. Here $ba$ denotes categorical composition. The linearization has the nonzero rays as a basis and extends this product bilinearly. An algebra is \emph{ray-standard} when it is isomorphic to this linearization. A category is \emph{mild} when every proper ideal quotient is locally representation-finite.

A contour is a pair of paths with the same nonzero ray. It is essential if the equality cannot be obtained by replacing equal-ray proper subpaths, and non-deep if the common ray is not annihilated by both endpoint corner radicals. The full subcategory on its vertices is its support category. The structure and disjointness theorems in \cite[Theorem 6]{Bongartz13} identify the essential non-deep support categories as pennies, dumb-bells and diamonds, with distinct supports disjoint. The penny path coordinates needed here are specified in Section~\ref{sec:boundary}. In the stem category, morphisms are the subbimodules generated by rays and composition is their product. The mixed forms of \cite[Theorem 10]{Bongartz13} use stem composition on selected penny supports and ray composition elsewhere.

\begin{lemma}\label{lem:ray-mesh}
A ray-standard representation-finite algebra over an algebraically closed field is mesh-standard.
\end{lemma}
\begin{proof}
Let $\Gamma$ be its Auslander--Reiten quiver and $\mathcal C=k(\Gamma)$ its mesh category. By \cite[Theorem 14]{Bongartz13}, the universal mesh category covers the category of indecomposables. This covering preserves radicals. Indeed, if a noninvertible map upstairs had invertible image, lift an inverse using the direct-sum Hom isomorphism in the definition of a covering. The component returning to the original object gives a left inverse. Local endomorphism rings then make the resulting nonzero idempotent on the other object the identity, so the original map would be invertible. The contradiction proves the assertion.

The finite category of indecomposables has nilpotent radical, say of exponent $N$. Every path of length $N$ upstairs maps to zero, hence is zero, since the covering is injective on each individual Hom space. Every path in $\mathcal C$ lifts, so its length-$N$ paths vanish too. As $\Gamma$ is finite, $\mathcal C$ is finite-dimensional; the universal mesh map to $\mathcal C$ is a covering between locally bounded categories.

Apply \cite[Theorem 15]{Bongartz13} first to the covering from the universal mesh category to the indecomposable category of the given algebra. The universal mesh category is therefore an Auslander category. Apply the same theorem to its covering of $\mathcal C$; hence $\mathcal C$ is also an Auslander category and is the indecomposable category of its algebra of projective points. These points are detected by the projective dimensions of simple contravariant functors: the radical inclusion at a projective point gives a resolution of length at most one, whereas at a nonprojective point the minimal mesh resolution has a nonzero term in degree two. The projective points are therefore exactly those of $\Gamma$. By \cite[Theorem 18(c),(d)]{Bongartz13}, their algebra is the ray-standard form, which is the given algebra.
\end{proof}


